% ch5-valeurs-propres.tex — Chapitre 5 : Valeurs propres, vecteurs propres et diagonalisation
% Ce chapitre développe la théorie spectrale en dimension finie :
% valeurs propres, polynôme caractéristique, théorème de Cayley-Hamilton,
% diagonalisation, trigonalisation, polynôme minimal, et applications.

\chapter{Valeurs propres, vecteurs propres et diagonalisation}\label{ch:valeurs-propres}

\section*{Introduction}
\addcontentsline{toc}{section}{Introduction}

L'un des objectifs fondamentaux de l'algèbre linéaire est de comprendre
la structure d'un endomorphisme en le ramenant, si possible, à une forme
aussi simple que possible. La \emph{diagonalisation}\index{diagonalisation}
consiste à trouver une base dans laquelle la matrice de l'endomorphisme
est diagonale~: l'action se réduit alors à de simples multiplications
scalaires le long de directions privilégiées.

\medskip

Cette idée est au c\oe ur de nombreuses applications~:
\begin{itemize}
  \item \textbf{Équations différentielles linéaires~:}\index{équations différentielles}
    la résolution du système $X'(t)=AX(t)$ se ramène, lorsque $A$ est
    diagonalisable, à $n$ équations scalaires indépendantes
    $x_i'(t)=\lambda_i x_i(t)$.
  \item \textbf{Chaînes de Markov~:}\index{chaîne de Markov} le
    comportement asymptotique d'une chaîne de Markov est gouverné par
    les valeurs propres de sa matrice de transition.
  \item \textbf{Analyse vibratoire~:}\index{vibrations} les modes propres
    d'un système mécanique correspondent aux vecteurs propres de la
    matrice de raideur.
  \item \textbf{Analyse en composantes principales (ACP)~:}\index{ACP}
    la réduction de dimension repose sur les vecteurs propres de la
    matrice de covariance.
  \item \textbf{Calcul de puissances de matrices~:} si $A=P\diag(\lambda_1,
    \ldots,\lambda_n)\inv{P}$, alors $A^k=P\diag(\lambda_1^k,\ldots,
    \lambda_n^k)\inv{P}$, ce qui permet de résoudre des récurrences
    linéaires.
\end{itemize}

\medskip

Dans tout ce chapitre, $\K$ désigne un corps commutatif ($\R$ ou $\C$
en pratique), $E$ un $\K$-espace vectoriel de dimension finie
$n\geq 1$, et $f\in\End(E)$.

% ============================================================
\section{Valeurs propres et vecteurs propres}\label{sec:vp-def}
% ============================================================

\begin{definition}{Valeur propre, vecteur propre}{def:vp}
  Soit $f\in\End(E)$.\index{valeur propre}\index{vecteur propre}
  \begin{enumerate}[label=(\roman*)]
    \item Un scalaire $\lambda\in\K$ est une \emph{valeur propre} de $f$
      s'il existe un vecteur $\vec{v}\in E\setminus\set{\vzero}$ tel que
      \[
        f(\vec{v})=\lambda\vec{v}.
      \]
      Le vecteur $\vec{v}$ est alors appelé \emph{vecteur propre} de $f$
      associé à la valeur propre~$\lambda$.
    \item De manière équivalente, $\lambda$ est valeur propre de $f$ si
      et seulement si $f-\lambda\Id$ n'est pas injective, c'est-à-dire
      si et seulement si $\Ker(f-\lambda\Id)\neq\set{\vzero}$.
  \end{enumerate}
\end{definition}

\begin{definition}{Valeur propre et vecteur propre d'une matrice}{def:vp-mat}
  Soit $A\in\Mat_n(\K)$.\index{valeur propre!d'une matrice}
  Un scalaire $\lambda\in\K$ est une \emph{valeur propre} de $A$ s'il
  existe un vecteur colonne $X\in\K^n\setminus\set{\vzero}$ tel que
  $AX=\lambda X$, c'est-à-dire $(A-\lambda I_n)X=0$.
  Ceci revient à dire que $\det(A-\lambda I_n)=0$.
\end{definition}

\begin{remark}{Lien endomorphisme--matrice}{rem:vp-endo-mat}
  Si $\cB$ est une base de $E$ et $A=\Mat_\cB(f)$, alors $\lambda$ est
  valeur propre de $f$ si et seulement si $\lambda$ est valeur propre
  de~$A$. De plus, si $\vec{v}$ est vecteur propre de $f$ pour
  $\lambda$, alors le vecteur de coordonnées de $\vec{v}$ dans $\cB$
  est vecteur propre de $A$ pour $\lambda$.
\end{remark}

\begin{example}{Valeurs propres d'une matrice $2\times 2$}{ex:vp-2x2}
  Soit $A=\begin{pmatrix}3&1\\0&2\end{pmatrix}$. Cherchons les valeurs
  propres de~$A$.

  On résout $\det(A-\lambda I_2)=0$~:
  \[
    \det\begin{pmatrix}3-\lambda&1\\0&2-\lambda\end{pmatrix}
    =(3-\lambda)(2-\lambda)=0.
  \]
  Les valeurs propres sont $\lambda_1=3$ et $\lambda_2=2$.

  Pour $\lambda_1=3$~: $(A-3I_2)X=0$ donne
  $\begin{pmatrix}0&1\\0&-1\end{pmatrix}\begin{pmatrix}x\\y\end{pmatrix}
  =\begin{pmatrix}0\\0\end{pmatrix}$, soit $y=0$. Les vecteurs propres
  associés sont de la forme $\begin{pmatrix}x\\0\end{pmatrix}$, $x\neq 0$.

  Pour $\lambda_2=2$~: $(A-2I_2)X=0$ donne
  $\begin{pmatrix}1&1\\0&0\end{pmatrix}\begin{pmatrix}x\\y\end{pmatrix}
  =\begin{pmatrix}0\\0\end{pmatrix}$, soit $x=-y$. Les vecteurs propres
  associés sont de la forme $\begin{pmatrix}-y\\y\end{pmatrix}$, $y\neq 0$.
\end{example}

\begin{proposition}{Indépendance linéaire des vecteurs propres}{prop:vp-li}
  Soient $\lambda_1,\ldots,\lambda_r$ des valeurs propres
  \emph{distinctes} de $f$, et soient $\vec{v}_1,\ldots,\vec{v}_r$ des
  vecteurs propres associés respectivement. Alors la famille
  $(\vec{v}_1,\ldots,\vec{v}_r)$ est libre.
\end{proposition}

\begin{proof}
  Par récurrence sur~$r$.

  \emph{Initialisation~:} pour $r=1$, $\vec{v}_1\neq\vzero$ donc
  $(\vec{v}_1)$ est libre.

  \emph{Hérédité~:} supposons le résultat vrai pour $r-1$ vecteurs propres
  associés à des valeurs propres distinctes. Supposons que
  \begin{equation}\label{eq:comb-vp}
    \alpha_1\vec{v}_1+\alpha_2\vec{v}_2+\cdots+\alpha_r\vec{v}_r=\vzero.
  \end{equation}
  Appliquons $f$ à cette relation~:
  \[
    \alpha_1\lambda_1\vec{v}_1+\alpha_2\lambda_2\vec{v}_2+\cdots
    +\alpha_r\lambda_r\vec{v}_r=\vzero.
  \]
  En multipliant \eqref{eq:comb-vp} par $\lambda_r$ et en soustrayant~:
  \[
    \alpha_1(\lambda_1-\lambda_r)\vec{v}_1+\cdots
    +\alpha_{r-1}(\lambda_{r-1}-\lambda_r)\vec{v}_{r-1}=\vzero.
  \]
  Par hypothèse de récurrence, $(\vec{v}_1,\ldots,\vec{v}_{r-1})$ est
  libre. Comme les $\lambda_i-\lambda_r\neq 0$ pour $i<r$, on obtient
  $\alpha_i=0$ pour $i=1,\ldots,r-1$. En reportant dans
  \eqref{eq:comb-vp}~: $\alpha_r\vec{v}_r=\vzero$, et comme
  $\vec{v}_r\neq\vzero$, $\alpha_r=0$.
\end{proof}

\begin{corollary}{Nombre de valeurs propres}{cor:nb-vp}
  Un endomorphisme d'un espace de dimension $n$ possède au plus $n$
  valeurs propres distinctes.
\end{corollary}

% ============================================================
\section{Sous-espaces propres}\label{sec:sous-espaces-propres}
% ============================================================

\begin{definition}{Sous-espace propre}{def:sep}
  Soit $\lambda$ une valeur propre de $f\in\End(E)$.
  Le \emph{sous-espace propre}\index{sous-espace propre} associé à
  $\lambda$ est
  \[
    E_\lambda(f)\deff\Ker(f-\lambda\Id)
    =\setcond{\vec{v}\in E}{f(\vec{v})=\lambda\vec{v}}.
  \]
\end{definition}

\begin{proposition}{$E_\lambda(f)$ est un sous-espace vectoriel}{prop:sep-sev}
  Pour toute valeur propre $\lambda$ de $f$, $E_\lambda(f)$ est un
  sous-espace vectoriel de $E$, de dimension $\geq 1$.
\end{proposition}

\begin{proof}
  $E_\lambda(f)=\Ker(f-\lambda\Id)$ est le noyau d'un endomorphisme,
  donc un sous-espace vectoriel de~$E$. De plus, $\lambda$ étant valeur
  propre, il existe $\vec{v}\neq\vzero$ tel que $f(\vec{v})=\lambda\vec{v}$,
  donc $E_\lambda(f)\neq\set{\vzero}$ et $\dim E_\lambda(f)\geq 1$.
\end{proof}

\begin{proposition}{Somme directe des sous-espaces propres}{prop:somme-sep}
  Soient $\lambda_1,\ldots,\lambda_r$ les valeurs propres distinctes
  de~$f$. Alors la somme
  $E_{\lambda_1}(f)+E_{\lambda_2}(f)+\cdots+E_{\lambda_r}(f)$
  est directe. Autrement dit~:
  \[
    \sum_{i=1}^{r}E_{\lambda_i}(f)=\bigoplus_{i=1}^{r}E_{\lambda_i}(f).
  \]
\end{proposition}

\begin{proof}
  Il suffit de montrer que si $\vec{v}_1+\cdots+\vec{v}_r=\vzero$ avec
  $\vec{v}_i\in E_{\lambda_i}(f)$, alors $\vec{v}_i=\vzero$ pour tout~$i$.
  Parmi les $\vec{v}_i$ non nuls, chacun est vecteur propre pour
  $\lambda_i$. Par la 
ef{prop:prop:vp-li}, ces vecteurs propres (associés
  à des valeurs propres distinctes) forment une famille libre. La relation
  $\vec{v}_1+\cdots+\vec{v}_r=\vzero$ impose alors que tous les
  coefficients (ici tous égaux à~$1$) soient nuls, ce qui n'est possible
  que si tous les $\vec{v}_i$ non nuls n'existent pas, c'est-à-dire
  $\vec{v}_i=\vzero$ pour tout~$i$.
\end{proof}

\begin{example}{Sous-espaces propres d'une matrice $3\times 3$}{ex:sep-3x3}
  Soit $A=\begin{pmatrix}2&1&0\\0&2&0\\0&0&3\end{pmatrix}$.
  Les valeurs propres sont $\lambda_1=2$ et $\lambda_2=3$.

  Pour $\lambda_1=2$~:
  $A-2I_3=\begin{pmatrix}0&1&0\\0&0&0\\0&0&1\end{pmatrix}$,
  donc $E_2(A)=\Vect\left(\begin{pmatrix}1\\0\\0\end{pmatrix}\right)$.
  On a $\dim E_2(A)=1$.

  Pour $\lambda_2=3$~:
  $A-3I_3=\begin{pmatrix}-1&1&0\\0&-1&0\\0&0&0\end{pmatrix}$,
  donc $E_3(A)=\Vect\left(\begin{pmatrix}0\\0\\1\end{pmatrix}\right)$.
  On a $\dim E_3(A)=1$.

  Ici $\dim E_2(A)+\dim E_3(A)=2<3=\dim E$~: la matrice $A$ n'est
  \emph{pas} diagonalisable, comme nous le verrons plus loin.
\end{example}

% ============================================================
\section{Polynôme caractéristique}\label{sec:poly-car}
% ============================================================

\begin{definition}{Polynôme caractéristique}{def:poly-car}
  Soit $A\in\Mat_n(\K)$. Le \emph{polynôme caractéristique}%
  \index{polynôme caractéristique} de $A$ est le polynôme
  \[
    \chi_A(\lambda)\deff\det(A-\lambda I_n)\in\K[\lambda].
  \]
  Si $f\in\End(E)$ et $A=\Mat_\cB(f)$ pour une base $\cB$ de $E$, on
  pose $\chi_f\deff\chi_A$. Ce polynôme ne dépend pas du choix de la
  base~$\cB$.
\end{definition}

\begin{proof}[Indépendance du choix de la base]
  Si $\cB'$ est une autre base et $P$ la matrice de passage, alors
  $\Mat_{\cB'}(f)=\inv{P}AP$, et
  \[
    \det(\inv{P}AP-\lambda I_n)=\det(\inv{P}(A-\lambda I_n)P)
    =\det(\inv{P})\det(A-\lambda I_n)\det(P)=\det(A-\lambda I_n).
    \qedhere
  \]
\end{proof}

\begin{proposition}{Propriétés du polynôme caractéristique}{prop:poly-car-props}
  Soit $A\in\Mat_n(\K)$. Alors~:
  \begin{enumerate}[label=(\roman*)]
    \item $\chi_A$ est un polynôme de degré $n$ en $\lambda$, de
      coefficient dominant $(-1)^n$~:
      \[
        \chi_A(\lambda)=(-1)^n\lambda^n+(-1)^{n-1}\tr(A)\lambda^{n-1}
        +\cdots+\det(A).
      \]
    \item $\chi_A(0)=\det(A)$.
    \item Le coefficient de $(-1)^{n-1}\lambda^{n-1}$ est $\tr(A)$.
    \item $\lambda_0$ est valeur propre de $A$ si et seulement si
      $\chi_A(\lambda_0)=0$.
    \item Deux matrices semblables ont le même polynôme caractéristique.
  \end{enumerate}
\end{proposition}

\begin{proof}
  \emph{(i)} On développe $\det(A-\lambda I_n)$. Le terme de plus haut
  degré provient du produit diagonal
  $(a_{11}-\lambda)(a_{22}-\lambda)\cdots(a_{nn}-\lambda)$, ce qui donne
  $(-\lambda)^n+\tr(A)(-\lambda)^{n-1}+\cdots$, soit
  $(-1)^n\lambda^n+(-1)^{n-1}\tr(A)\lambda^{n-1}+\cdots$.

  \emph{(ii)} $\chi_A(0)=\det(A-0\cdot I_n)=\det(A)$.

  \emph{(iii)} Découle de (i).

  \emph{(iv)} $\lambda_0$ est valeur propre de $A$ ssi
  $\Ker(A-\lambda_0 I_n)\neq\set{\vzero}$ ssi $A-\lambda_0 I_n$ n'est
  pas inversible ssi $\det(A-\lambda_0 I_n)=0$ ssi
  $\chi_A(\lambda_0)=0$.

  \emph{(v)} Déjà prouvé dans la démonstration de l'indépendance du
  choix de base.
\end{proof}

\begin{example}{Polynôme caractéristique --- cas $2\times 2$}{ex:chi-2x2}
  Pour $A=\begin{pmatrix}a&b\\c&d\end{pmatrix}$~:
  \[
    \chi_A(\lambda)=\det\begin{pmatrix}a-\lambda&b\\c&d-\lambda\end{pmatrix}
    =(a-\lambda)(d-\lambda)-bc
    =\lambda^2-(a+d)\lambda+(ad-bc)
    =\lambda^2-\tr(A)\lambda+\det(A).
  \]
\end{example}

\begin{example}{Polynôme caractéristique --- cas $3\times 3$}{ex:chi-3x3}
  Soit $A=\begin{pmatrix}1&2&0\\0&3&1\\0&0&2\end{pmatrix}$
  (matrice triangulaire supérieure).
  Alors~:
  \[
    \chi_A(\lambda)=(1-\lambda)(3-\lambda)(2-\lambda)
    =-(\lambda-1)(\lambda-2)(\lambda-3).
  \]
  Les valeurs propres sont $1$, $2$ et $3$. On retrouve que le polynôme
  caractéristique d'une matrice triangulaire se lit directement sur la
  diagonale.
\end{example}

\begin{proposition}{Cas d'une matrice triangulaire}{prop:chi-triang}
  Si $A$ est triangulaire (supérieure ou inférieure) avec des
  coefficients diagonaux $a_{11},\ldots,a_{nn}$, alors~:
  \[
    \chi_A(\lambda)=\prod_{i=1}^n(a_{ii}-\lambda).
  \]
  En particulier, les valeurs propres de $A$ sont exactement les
  coefficients diagonaux.
\end{proposition}

% ============================================================
\section{Théorème de Cayley-Hamilton}\label{sec:cayley-hamilton}
% ============================================================

\begin{theorem}{Cayley-Hamilton}{thm:cayley-hamilton}
  \index{théorème de Cayley-Hamilton}
  Soit $A\in\Mat_n(\K)$ de polynôme caractéristique
  $\chi_A(\lambda)=(-1)^n\lambda^n+c_{n-1}\lambda^{n-1}+\cdots+c_1\lambda+c_0$.
  Alors~:
  \[
    \chi_A(A)=(-1)^n A^n+c_{n-1}A^{n-1}+\cdots+c_1 A+c_0 I_n=0.
  \]
  Autrement dit, toute matrice carrée est racine de son propre polynôme
  caractéristique.
\end{theorem}

\begin{proof}
  Soit $B(\lambda)\deff\,^{\mathrm{t}}\!\operatorname{com}(A-\lambda I_n)$
  la matrice adjointe (comatrice transposée) de $A-\lambda I_n$.
  Par la formule classique de l'inverse (ou de la comatrice)~:
  \[
    (A-\lambda I_n)\,B(\lambda)=\det(A-\lambda I_n)\cdot I_n
    =\chi_A(\lambda)\cdot I_n.
  \]
  Chaque coefficient de $B(\lambda)$ est un cofacteur de $A-\lambda I_n$,
  donc un polynôme en $\lambda$ de degré au plus $n-1$. On peut donc
  écrire~:
  \[
    B(\lambda)=B_0+B_1\lambda+B_2\lambda^2+\cdots+B_{n-1}\lambda^{n-1},
  \]
  où $B_0,B_1,\ldots,B_{n-1}\in\Mat_n(\K)$ sont des matrices constantes
  (ne dépendant pas de $\lambda$).

  D'autre part, $\chi_A(\lambda)=c_0+c_1\lambda+\cdots+c_{n-1}\lambda^{n-1}
  +(-1)^n\lambda^n$.

  En développant $(A-\lambda I_n)\,B(\lambda)$~:
  \[
    (A-\lambda I_n)(B_0+B_1\lambda+\cdots+B_{n-1}\lambda^{n-1})
    =\chi_A(\lambda)\cdot I_n.
  \]
  On développe le membre de gauche et on identifie les coefficients de
  chaque puissance de $\lambda$~:
  \begin{align*}
    \lambda^0 &: \quad AB_0=c_0 I_n, \\
    \lambda^1 &: \quad AB_1-B_0=c_1 I_n, \\
    \lambda^2 &: \quad AB_2-B_1=c_2 I_n, \\
    &\;\;\vdots \\
    \lambda^{k} &: \quad AB_k-B_{k-1}=c_k I_n, \\
    &\;\;\vdots \\
    \lambda^{n-1} &: \quad AB_{n-1}-B_{n-2}=c_{n-1} I_n, \\
    \lambda^n &: \quad -B_{n-1}=(-1)^n I_n.
  \end{align*}
  Multiplions la $k$-ième relation ($k=0,1,\ldots,n$) par $A^k$
  \emph{à gauche}~:
  \begin{align*}
    A^0\cdot AB_0 &= c_0 I_n, \\
    A^1(AB_1-B_0) &= c_1 A, \\
    A^2(AB_2-B_1) &= c_2 A^2, \\
    &\;\;\vdots \\
    A^{n-1}(AB_{n-1}-B_{n-2}) &= c_{n-1}A^{n-1}, \\
    A^n(-B_{n-1}) &= (-1)^n A^n.
  \end{align*}
  En sommant toutes ces relations, le membre de gauche se téléscope~:
  chaque terme $A^k B_{k-1}$ apparaît une fois avec le signe~$+$ (dans
  la ligne $k$) et une fois avec le signe~$-$ (dans la ligne $k-1$
  sous la forme $A^{k-1}\cdot AB_{k-1}=A^k B_{k-1}$). Il reste~:
  \[
    0=c_0 I_n+c_1 A+c_2 A^2+\cdots+c_{n-1}A^{n-1}+(-1)^n A^n
    =\chi_A(A).
    \qedhere
  \]
\end{proof}

\begin{example}{Vérification de Cayley-Hamilton}{ex:cayley-hamilton}
  Soit $A=\begin{pmatrix}1&2\\3&4\end{pmatrix}$. On a
  $\chi_A(\lambda)=\lambda^2-5\lambda-2$. Vérifions~:
  \[
    A^2-5A-2I_2=\begin{pmatrix}7&10\\15&22\end{pmatrix}
    -\begin{pmatrix}5&10\\15&20\end{pmatrix}
    -\begin{pmatrix}2&0\\0&2\end{pmatrix}
    =\begin{pmatrix}0&0\\0&0\end{pmatrix}. \quad\checkmark
  \]
\end{example}

% ============================================================
\section{Spectre d'un endomorphisme}\label{sec:spectre}
% ============================================================

\begin{definition}{Spectre}{def:spectre}
  \index{spectre}
  Le \emph{spectre} de $f\in\End(E)$ est l'ensemble de ses valeurs
  propres~:
  \[
    \Spec(f)\deff\setcond{\lambda\in\K}{\chi_f(\lambda)=0}.
  \]
  De même, le spectre d'une matrice $A\in\Mat_n(\K)$ est
  $\Spec(A)\deff\setcond{\lambda\in\K}{\chi_A(\lambda)=0}$.
\end{definition}

\begin{remark}{Dépendance en le corps}{rem:spectre-corps}
  Le spectre dépend du corps~$\K$. Par exemple, la matrice de rotation
  $R=\begin{pmatrix}0&-1\\1&0\end{pmatrix}$ a pour polynôme
  caractéristique $\chi_R(\lambda)=\lambda^2+1$. Si $\K=\R$, alors
  $\Spec(R)=\emptyset$~; si $\K=\C$, alors $\Spec(R)=\set{i,-i}$.
\end{remark}

% ============================================================
\section{Multiplicité algébrique et multiplicité géométrique}\label{sec:multiplicites}
% ============================================================

\begin{definition}{Multiplicité algébrique}{def:mult-alg}
  \index{multiplicité algébrique}
  Soit $\lambda$ une valeur propre de $f$. La \emph{multiplicité
  algébrique} de $\lambda$, notée $m_a(\lambda)$, est l'ordre de
  multiplicité de $\lambda$ comme racine de $\chi_f$. Autrement dit,
  $m_a(\lambda)$ est le plus grand entier $k$ tel que
  $(\lambda_0-\lambda)^k$ divise $\chi_f(\lambda)$ (où $\lambda_0$
  est la variable du polynôme), c'est-à-dire~:
  \[
    \chi_f(\lambda_0)=(\lambda_0-\lambda)^{m_a(\lambda)}\,Q(\lambda_0),
    \quad Q(\lambda)\neq 0.
  \]
\end{definition}

\begin{definition}{Multiplicité géométrique}{def:mult-geo}
  \index{multiplicité géométrique}
  La \emph{multiplicité géométrique} de la valeur propre $\lambda$
  est la dimension du sous-espace propre associé~:
  \[
    m_g(\lambda)\deff\dim E_\lambda(f)=\dim\Ker(f-\lambda\Id).
  \]
\end{definition}

\begin{theorem}{Encadrement des multiplicités}{thm:encadrement-mult}
  Pour toute valeur propre $\lambda$ de $f$~:
  \[
    1\leq m_g(\lambda)\leq m_a(\lambda).
  \]
\end{theorem}

\begin{proof}
  L'inégalité $m_g(\lambda)\geq 1$ résulte de la définition~: $\lambda$
  étant valeur propre, $E_\lambda(f)\neq\set{\vzero}$ donc
  $\dim E_\lambda(f)\geq 1$.

  Montrons que $m_g(\lambda)\leq m_a(\lambda)$. Posons $p=m_g(\lambda)$
  et choisissons une base $(\vec{v}_1,\ldots,\vec{v}_p)$ de
  $E_\lambda(f)$. Complétons-la en une base
  $\cB=(\vec{v}_1,\ldots,\vec{v}_p,\vec{w}_{p+1},\ldots,\vec{w}_n)$
  de~$E$. Dans cette base, la matrice de $f$ s'écrit~:
  \[
    \Mat_\cB(f)=\begin{pmatrix}\lambda I_p & C \\ 0 & D\end{pmatrix}
  \]
  où $C\in\Mat_{p,n-p}(\K)$ et $D\in\Mat_{n-p}(\K)$. Alors~:
  \[
    \chi_f(\mu)=\det\begin{pmatrix}(\lambda-\mu)I_p & C \\ 0 & D-\mu I_{n-p}
    \end{pmatrix}=(\lambda-\mu)^p\det(D-\mu I_{n-p}).
  \]
  Donc $(\lambda-\mu)^p$ divise $\chi_f(\mu)$, ce qui donne
  $m_a(\lambda)\geq p=m_g(\lambda)$.
\end{proof}

\begin{example}{Multiplicités distinctes}{ex:mult-distinctes}
  Pour $A=\begin{pmatrix}2&1&0\\0&2&0\\0&0&3\end{pmatrix}$ (cf.\
  
ef{ex:ex:sep-3x3}), on a $\chi_A(\lambda)=-(2-\lambda)^2(3-\lambda)$.
  \begin{itemize}
    \item Pour $\lambda=2$~: $m_a(2)=2$ et $m_g(2)=\dim E_2(A)=1$.
      Ici $m_g<m_a$.
    \item Pour $\lambda=3$~: $m_a(3)=1$ et $m_g(3)=\dim E_3(A)=1$.
      Ici $m_g=m_a$.
  \end{itemize}
\end{example}

% ============================================================
\section{Diagonalisabilité}\label{sec:diagonalisabilite}
% ============================================================

\begin{definition}{Endomorphisme diagonalisable}{def:diag}
  \index{diagonalisable!endomorphisme}
  Un endomorphisme $f\in\End(E)$ est dit \emph{diagonalisable} s'il
  existe une base de $E$ formée de vecteurs propres de~$f$. De manière
  équivalente, $f$ est diagonalisable s'il existe une base $\cB$ de $E$
  telle que $\Mat_\cB(f)$ soit une matrice diagonale.
\end{definition}

\begin{definition}{Matrice diagonalisable}{def:mat-diag}
  \index{diagonalisable!matrice}
  Une matrice $A\in\Mat_n(\K)$ est dite \emph{diagonalisable} s'il
  existe $P\in\GL_n(\K)$ telle que $\inv{P}AP$ soit diagonale.
  Autrement dit, $A$ est semblable à une matrice diagonale.
\end{definition}

\begin{remark}{Diagonalisation concrète}{rem:diag-concrete}
  Si $f$ est diagonalisable et si
  $\cB=(\vec{v}_1,\ldots,\vec{v}_n)$ est une base de vecteurs propres
  avec $f(\vec{v}_i)=\lambda_i\vec{v}_i$, alors~:
  \[
    \Mat_\cB(f)=\diag(\lambda_1,\ldots,\lambda_n)
    =\begin{pmatrix}\lambda_1&&0\\&\ddots&\\0&&\lambda_n\end{pmatrix}.
  \]
  Si $A=\Mat_\cE(f)$ dans la base canonique et $P$ est la matrice dont
  les colonnes sont les coordonnées des $\vec{v}_i$, alors
  $A=P\diag(\lambda_1,\ldots,\lambda_n)\inv{P}$.
\end{remark}

\begin{theorem}{Critère de diagonalisabilité}{thm:critere-diag}
  \index{diagonalisation!critère}
  Soit $f\in\End(E)$ avec $\dim E=n$. Notons
  $\lambda_1,\ldots,\lambda_r$ les valeurs propres distinctes de~$f$.
  Les assertions suivantes sont équivalentes~:
  \begin{enumerate}[label=(\roman*)]
    \item $f$ est diagonalisable.
    \item $E=\bigoplus_{i=1}^{r}E_{\lambda_i}(f)$.
    \item $\sum_{i=1}^{r}\dim E_{\lambda_i}(f)=n$.
    \item Le polynôme caractéristique $\chi_f$ est scindé sur $\K$
      \emph{et} $m_g(\lambda_i)=m_a(\lambda_i)$ pour tout
      $i\in\set{1,\ldots,r}$.
  \end{enumerate}
\end{theorem}

\begin{proof}
  \emph{(i) $\Rightarrow$ (ii)~:} Si $f$ est diagonalisable, il existe
  une base $\cB=(\vec{v}_1,\ldots,\vec{v}_n)$ de vecteurs propres.
  Chaque $\vec{v}_j$ appartient à un sous-espace propre
  $E_{\lambda_{i(j)}}(f)$. Comme $\cB$ engendre $E$, on a
  $E=E_{\lambda_1}(f)+\cdots+E_{\lambda_r}(f)$. La somme est directe
  par la 
ef{prop:prop:somme-sep}.

  \emph{(ii) $\Rightarrow$ (iii)~:} Immédiat~: si la somme directe
  vaut $E$, la somme des dimensions vaut $n$.

  \emph{(iii) $\Rightarrow$ (iv)~:} On a $\sum_i m_g(\lambda_i)=n$.
  Or $\sum_i m_a(\lambda_i)\leq n$ (car $\deg\chi_f=n$) et
  $m_g(\lambda_i)\leq m_a(\lambda_i)$ par le
  
ef{thm:thm:encadrement-mult}. Donc~:
  \[
    n=\sum_{i=1}^r m_g(\lambda_i)\leq\sum_{i=1}^r m_a(\lambda_i)\leq n.
  \]
  Toutes les inégalités sont des égalités. En particulier,
  $\sum_i m_a(\lambda_i)=n$, ce qui signifie que $\chi_f$ est scindé
  sur~$\K$ (toutes les racines sont dans~$\K$ et leurs multiplicités
  totalisent $n=\deg\chi_f$). De plus,
  $m_g(\lambda_i)=m_a(\lambda_i)$ pour tout~$i$.

  \emph{(iv) $\Rightarrow$ (i)~:} Si $\chi_f$ est scindé et
  $m_g(\lambda_i)=m_a(\lambda_i)$ pour tout~$i$, alors
  \[
    \sum_{i=1}^r\dim E_{\lambda_i}(f)=\sum_{i=1}^r m_g(\lambda_i)
    =\sum_{i=1}^r m_a(\lambda_i)=n.
  \]
  Par la 
ef{prop:prop:somme-sep}, la somme des sous-espaces propres est
  directe et de dimension~$n$, donc $E=\bigoplus_i E_{\lambda_i}(f)$.
  En prenant la réunion de bases de chaque $E_{\lambda_i}(f)$, on
  obtient une base de $E$ formée de vecteurs propres.
\end{proof}

\begin{corollary}{$n$ valeurs propres distinctes}{cor:n-vp-diag}
  Si $f$ possède $n=\dim E$ valeurs propres distinctes, alors $f$ est
  diagonalisable.
\end{corollary}

\begin{proof}
  Si les $n$ valeurs propres sont distinctes, chacune a multiplicité
  algébrique~$1$, donc multiplicité géométrique~$1$ (par l'encadrement).
  Le critère (iv) du 
ef{thm:thm:critere-diag} est vérifié.
\end{proof}

% ============================================================
\section{Algorithme de diagonalisation}\label{sec:algo-diag}
% ============================================================

Voici la procédure à suivre pour diagonaliser un endomorphisme $f$
(ou une matrice $A$)~:

\begin{enumerate}[label=\textbf{Étape \arabic*.},leftmargin=3em]
  \item \textbf{Calculer le polynôme caractéristique} $\chi_f(\lambda)
    =\det(A-\lambda I_n)$.
  \item \textbf{Trouver les racines} de $\chi_f$ dans $\K$. Ce sont
    les valeurs propres $\lambda_1,\ldots,\lambda_r$ avec leurs
    multiplicités algébriques $m_a(\lambda_i)$.
  \item \textbf{Vérifier que $\chi_f$ est scindé}~: la somme des
    multiplicités algébriques doit valoir~$n$. Sinon, $f$ n'est pas
    diagonalisable (sur~$\K$).
  \item \textbf{Calculer les sous-espaces propres}
    $E_{\lambda_i}=\Ker(A-\lambda_i I_n)$ en résolvant les systèmes
    linéaires.
  \item \textbf{Vérifier} que $\dim E_{\lambda_i}=m_a(\lambda_i)$ pour
    tout~$i$. Si c'est le cas, $f$ est diagonalisable~; sinon, non.
  \item \textbf{Former la matrice de passage} $P$ en juxtaposant les
    bases des sous-espaces propres. On obtient
    $\inv{P}AP=\diag(\lambda_1,\ldots,\lambda_1,\lambda_2,\ldots,
    \lambda_r)$.
\end{enumerate}

\begin{example}{Diagonalisation complète}{ex:diag-complete}
  Diagonalisons la matrice
  $A=\begin{pmatrix}4&1\\2&3\end{pmatrix}$.

  \textbf{Étape 1.} $\chi_A(\lambda)=\det\begin{pmatrix}4-\lambda&1\\
  2&3-\lambda\end{pmatrix}=(4-\lambda)(3-\lambda)-2
  =\lambda^2-7\lambda+10=(\lambda-5)(\lambda-2)$.

  \textbf{Étape 2.} Les valeurs propres sont $\lambda_1=5$ et
  $\lambda_2=2$, chacune de multiplicité algébrique~$1$.

  \textbf{Étape 3.} $\chi_A$ est scindé sur $\R$.

  \textbf{Étape 4.}
  \begin{itemize}
    \item $E_5$~: $(A-5I_2)X=0$, soit
      $\begin{pmatrix}-1&1\\2&-2\end{pmatrix}X=0$, d'où $x=y$.
      Donc $E_5=\Vect\begin{pmatrix}1\\1\end{pmatrix}$.
    \item $E_2$~: $(A-2I_2)X=0$, soit
      $\begin{pmatrix}2&1\\2&1\end{pmatrix}X=0$, d'où $2x+y=0$.
      Donc $E_2=\Vect\begin{pmatrix}1\\-2\end{pmatrix}$.
  \end{itemize}

  \textbf{Étape 5.} $\dim E_5=1=m_a(5)$ et $\dim E_2=1=m_a(2)$.
  La matrice est diagonalisable.

  \textbf{Étape 6.}
  $P=\begin{pmatrix}1&1\\1&-2\end{pmatrix}$,
  $\inv{P}=\frac{1}{-3}\begin{pmatrix}-2&-1\\-1&1\end{pmatrix}$, et
  \[
    \inv{P}AP=\begin{pmatrix}5&0\\0&2\end{pmatrix}.
  \]
\end{example}

\begin{example}{Diagonalisation d'une matrice $3\times 3$}{ex:diag-3x3}
  Soit $A=\begin{pmatrix}1&1&0\\0&2&0\\0&1&3\end{pmatrix}$.

  \textbf{Étape 1.}
  $\chi_A(\lambda)=\det\begin{pmatrix}1-\lambda&1&0\\0&2-\lambda&0\\
  0&1&3-\lambda\end{pmatrix}
  =(1-\lambda)(2-\lambda)(3-\lambda)$.

  \textbf{Étape 2.} Les valeurs propres sont $\lambda_1=1$,
  $\lambda_2=2$, $\lambda_3=3$, chacune simple.

  \textbf{Étape 3.} $\chi_A$ est scindé sur $\R$ et les trois valeurs
  propres sont distinctes~: $A$ est diagonalisable
  (
ef{cor:cor:n-vp-diag}).

  \textbf{Étape 4.}
  \begin{itemize}
    \item $E_1$~: $(A-I_3)X=0$ avec
      $A-I_3=\begin{pmatrix}0&1&0\\0&1&0\\0&1&2\end{pmatrix}$, d'où
      $y=0$, $z=0$. Donc
      $E_1=\Vect\begin{pmatrix}1\\0\\0\end{pmatrix}$.
    \item $E_2$~: $(A-2I_3)X=0$ avec
      $A-2I_3=\begin{pmatrix}-1&1&0\\0&0&0\\0&1&1\end{pmatrix}$, d'où
      $x=y$ et $y=-z$. Donc
      $E_2=\Vect\begin{pmatrix}1\\1\\-1\end{pmatrix}$.
    \item $E_3$~: $(A-3I_3)X=0$ avec
      $A-3I_3=\begin{pmatrix}-2&1&0\\0&-1&0\\0&1&0\end{pmatrix}$, d'où
      $y=0$ et $x=0$. Donc
      $E_3=\Vect\begin{pmatrix}0\\0\\1\end{pmatrix}$.

      Vérifions~: $A\begin{pmatrix}0\\0\\1\end{pmatrix}
      =\begin{pmatrix}0\\0\\3\end{pmatrix}
      =3\begin{pmatrix}0\\0\\1\end{pmatrix}$. \checkmark
  \end{itemize}

  \textbf{Étape 6.}
  $P=\begin{pmatrix}1&1&0\\0&1&0\\0&-1&1\end{pmatrix}$ et
  $\inv{P}AP=\diag(1,2,3)$.
\end{example}

% ============================================================
\section{Trigonalisation}\label{sec:trigonalisation}
% ============================================================

Lorsqu'un endomorphisme n'est pas diagonalisable, on peut chercher la
meilleure forme simple suivante~: une matrice triangulaire.

\begin{definition}{Endomorphisme trigonalisable}{def:trigonalisable}
  \index{trigonalisable}
  Un endomorphisme $f\in\End(E)$ est dit \emph{trigonalisable} s'il
  existe une base de $E$ dans laquelle la matrice de $f$ est
  triangulaire supérieure. Une matrice $A$ est trigonalisable si elle
  est semblable à une matrice triangulaire supérieure.
\end{definition}

\begin{theorem}{Critère de trigonalisation}{thm:trigonalisation}
  \index{trigonalisation}
  Un endomorphisme $f\in\End(E)$ est trigonalisable si et seulement si
  son polynôme caractéristique $\chi_f$ est scindé sur~$\K$.

  En particulier, sur $\K=\C$, \textbf{toute} matrice est
  trigonalisable (car tout polynôme est scindé sur $\C$ d'après le
  théorème fondamental de l'algèbre).
\end{theorem}

\begin{proof}
  \emph{$(\Rightarrow)$} Si $f$ est trigonalisable, il existe une base
  dans laquelle sa matrice est triangulaire supérieure avec les
  coefficients diagonaux $t_{11},\ldots,t_{nn}$. Par la
  
ef{prop:prop:chi-triang}, $\chi_f(\lambda)=\prod_{i=1}^n(t_{ii}-\lambda)$,
  qui est scindé sur~$\K$.

  \emph{$(\Leftarrow)$} Par récurrence sur $n=\dim E$.

  \emph{Initialisation~:} si $n=1$, toute matrice $1\times 1$ est
  triangulaire.

  \emph{Hérédité~:} soit $n\geq 2$. Comme $\chi_f$ est scindé, $f$
  admet au moins une valeur propre $\lambda_1\in\K$. Soit
  $\vec{v}_1\neq\vzero$ un vecteur propre associé. Soit $H$ un
  supplémentaire de $\Vect(\vec{v}_1)$ dans $E$ (un sous-espace tel
  que $E=\Vect(\vec{v}_1)\oplus H$, de dimension $n-1$). Posons
  $\cB_1=(\vec{v}_1,\vec{h}_2,\ldots,\vec{h}_n)$ une base de $E$
  adaptée à cette décomposition. Dans cette base~:
  \[
    \Mat_{\cB_1}(f)=\begin{pmatrix}\lambda_1 & * \\ 0 & A'\end{pmatrix}
  \]
  avec $A'\in\Mat_{n-1}(\K)$. Alors
  $\chi_f(\lambda)=(\lambda_1-\lambda)\chi_{A'}(\lambda)$, et comme
  $\chi_f$ est scindé, $\chi_{A'}$ l'est aussi. Par hypothèse de
  récurrence, $A'$ est trigonalisable~: il existe $Q\in\GL_{n-1}(\K)$
  telle que $\inv{Q}A'Q$ soit triangulaire supérieure. En posant
  $\tilde{P}=\begin{pmatrix}1&0\\0&Q\end{pmatrix}$, on obtient une
  matrice triangulaire supérieure pour $f$ dans la base correspondante.
\end{proof}

\begin{corollary}{Trigonalisation sur $\C$}{cor:trig-C}
  Toute matrice $A\in\Mat_n(\C)$ est trigonalisable~: il existe
  $P\in\GL_n(\C)$ et une matrice triangulaire supérieure~$T$ telles que
  $A=PTP^{-1}$. Les coefficients diagonaux de $T$ sont les valeurs
  propres de~$A$ (comptées avec multiplicité).
\end{corollary}

% ============================================================
\section{Polynôme minimal}\label{sec:poly-min}
% ============================================================

\begin{definition}{Polynôme minimal}{def:poly-min}
  \index{polynôme minimal}
  Soit $f\in\End(E)$. Le \emph{polynôme minimal} de $f$, noté
  $\mu_f$, est le polynôme unitaire de plus petit degré tel que
  $\mu_f(f)=0$ (l'endomorphisme nul).

  De même, le polynôme minimal de $A\in\Mat_n(\K)$ est le polynôme
  unitaire $\mu_A$ de plus petit degré tel que $\mu_A(A)=0$.
\end{definition}

\begin{proposition}{Existence et unicité}{prop:poly-min-existe}
  Le polynôme minimal existe et est unique. De plus, un polynôme
  $P\in\K[X]$ vérifie $P(f)=0$ si et seulement si $\mu_f$ divise~$P$.
\end{proposition}

\begin{proof}
  L'espace $\End(E)$ est de dimension $n^2$, donc la famille
  $(\Id,f,f^2,\ldots,f^{n^2})$ de $n^2+1$ éléments est liée~: il
  existe une relation non triviale $\sum_{k=0}^{n^2}a_k f^k=0$. Ceci
  montre l'existence d'un polynôme annulateur non nul (et le théorème
  de Cayley-Hamilton en fournit un de degré~$n$). Parmi tous les
  polynômes annulateurs non nuls, prenons-en un de degré minimal et
  rendons-le unitaire~: c'est $\mu_f$.

  \emph{Unicité~:} si $\mu'$ est un autre polynôme unitaire annulateur
  de degré minimal, alors $\mu_f-\mu'$ est annulateur de degré
  strictement inférieur à $\deg\mu_f$, donc $\mu_f-\mu'=0$.

  \emph{Divisibilité~:} si $P(f)=0$, effectuons la division euclidienne
  $P=Q\mu_f+R$ avec $\deg R<\deg\mu_f$. Alors $R(f)=P(f)-Q(f)\mu_f(f)
  =0-0=0$. Par minimalité de $\deg\mu_f$, on a $R=0$, donc
  $\mu_f\mid P$.
\end{proof}

\begin{theorem}{Relations entre $\mu_f$ et $\chi_f$}{thm:mu-chi}
  \begin{enumerate}[label=(\roman*)]
    \item $\mu_f$ divise $\chi_f$ (conséquence du théorème de
      Cayley-Hamilton).
    \item $\mu_f$ et $\chi_f$ ont les mêmes racines (dans une clôture
      algébrique de~$\K$).
    \item $\deg\mu_f\leq n=\deg\chi_f$.
  \end{enumerate}
\end{theorem}

\begin{proof}
  \emph{(i)} Par le 
ef{thm:thm:cayley-hamilton}, $\chi_f(f)=0$, donc
  $\mu_f\mid\chi_f$ par la 
ef{prop:prop:poly-min-existe}.

  \emph{(ii)} Toute racine de $\mu_f$ est racine de $\chi_f$ (car
  $\mu_f\mid\chi_f$). Réciproquement, soit $\lambda$ une racine de
  $\chi_f$~: c'est une valeur propre de $f$. Soit $\vec{v}\neq\vzero$
  un vecteur propre associé. Alors $\mu_f(f)\vec{v}=\mu_f(\lambda)\vec{v}
  =\vzero$. Comme $\vec{v}\neq\vzero$, on a $\mu_f(\lambda)=0$.

  \emph{(iii)} Conséquence immédiate de (i).
\end{proof}

\begin{theorem}{Polynôme minimal et diagonalisabilité}{thm:mu-diag}
  \index{diagonalisation!et polynôme minimal}
  L'endomorphisme $f$ est diagonalisable si et seulement si $\mu_f$
  est scindé à racines simples sur~$\K$, c'est-à-dire~:
  \[
    \mu_f(\lambda)=(\lambda-\lambda_1)(\lambda-\lambda_2)\cdots
    (\lambda-\lambda_r)
  \]
  où $\lambda_1,\ldots,\lambda_r$ sont les valeurs propres distinctes
  de~$f$.
\end{theorem}

\begin{proof}
  \emph{$(\Rightarrow)$} Supposons $f$ diagonalisable. Alors il existe
  une base dans laquelle
  $\Mat(f)=\diag(\lambda_1,\ldots,\lambda_1,\lambda_2,\ldots,\lambda_r)$.
  Posons $P(\lambda)=(\lambda-\lambda_1)\cdots(\lambda-\lambda_r)$
  (racines simples). On a
  $P(f)=(f-\lambda_1\Id)\circ\cdots\circ(f-\lambda_r\Id)$.
  Dans la base de diagonalisation, chaque vecteur de base $\vec{v}_j$
  est vecteur propre pour une certaine valeur propre $\lambda_{i_j}$,
  donc $(f-\lambda_{i_j}\Id)\vec{v}_j=\vzero$, et par conséquent
  $P(f)\vec{v}_j=\vzero$. Comme ceci vaut pour tous les vecteurs de
  la base, $P(f)=0$. Donc $\mu_f\mid P$. Comme $\mu_f$ et $\chi_f$
  ont les mêmes racines et que $P$ a les mêmes racines que $\mu_f$,
  le polynôme $\mu_f$ est à racines simples. Comme $\mu_f\mid P$ et
  $P$ est à racines simples, $\mu_f$ est aussi à racines simples.

  \emph{$(\Leftarrow)$} Supposons
  $\mu_f=(\lambda-\lambda_1)\cdots(\lambda-\lambda_r)$ scindé à racines
  simples. Alors $\mu_f(f)=(f-\lambda_1\Id)\circ\cdots\circ
  (f-\lambda_r\Id)=0$. Par le lemme des noyaux (appliqué aux polynômes
  deux à deux premiers entre eux $\lambda-\lambda_i$)~:
  \[
    E=\bigoplus_{i=1}^r\Ker(f-\lambda_i\Id)
    =\bigoplus_{i=1}^r E_{\lambda_i}(f).
  \]
  Donc $f$ est diagonalisable.
\end{proof}

\begin{example}{Polynôme minimal}{ex:poly-min}
  Soit $A=\begin{pmatrix}2&0&0\\0&2&0\\0&0&3\end{pmatrix}$. On a
  $\chi_A(\lambda)=-(\lambda-2)^2(\lambda-3)$.

  \begin{itemize}
    \item $E_2(A)=\Vect\left(\begin{pmatrix}1\\0\\0\end{pmatrix},
      \begin{pmatrix}0\\1\\0\end{pmatrix}\right)$, de dimension $2$.
    \item $E_3(A)=\Vect\left(\begin{pmatrix}0\\0\\1\end{pmatrix}\right)$.
  \end{itemize}
  La matrice $A$ est diagonalisable (c'est déjà une matrice diagonale).
  Le polynôme minimal est $\mu_A(\lambda)=(\lambda-2)(\lambda-3)$
  (racines simples), de degré~$2<3=\deg\chi_A$.
\end{example}

% ============================================================
\section{Applications}\label{sec:applications-vp}
% ============================================================

\subsection{Calcul de $A^k$}\label{subsec:puissances}

\begin{proposition}{Puissances d'une matrice diagonalisable}{prop:puissances}
  Si $A=P\,\diag(\lambda_1,\ldots,\lambda_n)\,\inv{P}$, alors pour
  tout $k\in\N$~:
  \[
    A^k=P\,\diag(\lambda_1^k,\ldots,\lambda_n^k)\,\inv{P}.
  \]
\end{proposition}

\begin{proof}
  Par récurrence immédiate~:
  $A^k=(P D\inv{P})^k=P D^k\inv{P}=P\diag(\lambda_1^k,\ldots,
  \lambda_n^k)\inv{P}$.
\end{proof}

\subsection{Récurrences linéaires~: la suite de Fibonacci}\label{subsec:fibonacci}
\index{suite de Fibonacci}

La suite de Fibonacci est définie par $F_0=0$, $F_1=1$ et
$F_{n+2}=F_{n+1}+F_n$. Posons
$X_n=\begin{pmatrix}F_{n+1}\\F_n\end{pmatrix}$. Alors~:
\[
  X_n=AX_{n-1}\quad\text{avec}\quad
  A=\begin{pmatrix}1&1\\1&0\end{pmatrix}.
\]
Donc $X_n=A^n X_0$ avec $X_0=\begin{pmatrix}1\\0\end{pmatrix}$.

Diagonalisons~$A$~: $\chi_A(\lambda)=\lambda^2-\lambda-1$, de racines
$\vphi=\frac{1+\sqrt{5}}{2}$ (nombre d'or\index{nombre d'or}) et
$\hat{\vphi}=\frac{1-\sqrt{5}}{2}$.

Les vecteurs propres sont
$\vec{v}_1=\begin{pmatrix}\vphi\\1\end{pmatrix}$ et
$\vec{v}_2=\begin{pmatrix}\hat{\vphi}\\1\end{pmatrix}$.
Avec $P=\begin{pmatrix}\vphi&\hat{\vphi}\\1&1\end{pmatrix}$, on a
$A=P\diag(\vphi,\hat{\vphi})\inv{P}$, d'où~:
\[
  A^n=P\begin{pmatrix}\vphi^n&0\\0&\hat{\vphi}^n\end{pmatrix}\inv{P}.
\]
En calculant la deuxième composante de $A^n X_0$, on obtient la
\textbf{formule de Binet}\index{formule de Binet}~:
\[
  \boxed{F_n=\frac{\vphi^n-\hat{\vphi}^n}{\sqrt{5}}
  =\frac{1}{\sqrt{5}}\left[\left(\frac{1+\sqrt{5}}{2}\right)^{\!n}
  -\left(\frac{1-\sqrt{5}}{2}\right)^{\!n}\right].}
\]

\subsection{Chaînes de Markov~: état stationnaire}\label{subsec:markov-ch5}
\index{chaîne de Markov}

Considérons une chaîne de Markov à deux états avec matrice de transition
\[
  M=\begin{pmatrix}0{,}7&0{,}4\\0{,}3&0{,}6\end{pmatrix}.
\]
(Les colonnes somment à $1$~: matrice stochastique.)

On a $\chi_M(\lambda)=(\lambda-1)(\lambda-0{,}3)$, de valeurs propres
$\lambda_1=1$ et $\lambda_2=0{,}3$.

Le vecteur propre associé à $\lambda_1=1$ (vecteur de probabilités
stationnaires) vérifie $M\pi=\pi$, soit
$\pi=\frac{1}{7}\begin{pmatrix}4\\3\end{pmatrix}$ (normalisé pour
que la somme des composantes vaille~$1$).

Comme $|\lambda_2|=0{,}3<1$, pour toute distribution initiale
$X_0$, la suite $X_n=M^n X_0$ converge vers~$\pi$~:
l'état stationnaire est atteint.

\subsection{Systèmes d'équations différentielles linéaires}\label{subsec:edo}
\index{équations différentielles!systèmes}

Considérons le système $X'(t)=AX(t)$ avec
$A=\begin{pmatrix}-1&1\\0&-2\end{pmatrix}$.

On a $\chi_A(\lambda)=(\lambda+1)(\lambda+2)$, de valeurs propres
$\lambda_1=-1$ et $\lambda_2=-2$.
\begin{itemize}
  \item $E_{-1}$~: vecteur propre
    $\vec{v}_1=\begin{pmatrix}1\\0\end{pmatrix}$.
  \item $E_{-2}$~: $(A+2I)X=0$ donne
    $\begin{pmatrix}1&1\\0&0\end{pmatrix}X=0$, soit
    $\vec{v}_2=\begin{pmatrix}-1\\1\end{pmatrix}$.
\end{itemize}

Dans la base $(\vec{v}_1,\vec{v}_2)$, le système devient
$Y'(t)=\diag(-1,-2)Y(t)$, de solution $y_1(t)=c_1 e^{-t}$,
$y_2(t)=c_2 e^{-2t}$. En revenant à la base canonique~:
\[
  X(t)=c_1 e^{-t}\begin{pmatrix}1\\0\end{pmatrix}
  +c_2 e^{-2t}\begin{pmatrix}-1\\1\end{pmatrix}
  =\begin{pmatrix}c_1 e^{-t}-c_2 e^{-2t}\\c_2 e^{-2t}\end{pmatrix}.
\]

% ============================================================
\section{Illustrations}\label{sec:illustrations-vp}
% ============================================================

\begin{figure}[ht]
\centering
\begin{tikzpicture}[scale=1.4, >=Stealth]
  % Axes
  \draw[gray!30, very thin] (-2.5,-2.5) grid (3.5,3.5);
  \draw[->] (-2.5,0) -- (3.7,0) node[right] {$x$};
  \draw[->] (0,-2.5) -- (0,3.7) node[above] {$y$};

  % Eigenvector v1 (lambda=2, stretching)
  \draw[->, thick, blue!70!black] (0,0) -- (1,1)
    node[above left] {$\vec{v}_1$};
  \draw[->, thick, blue, dashed] (0,0) -- (2,2)
    node[above left] {$f(\vec{v}_1)=2\vec{v}_1$};

  % Eigenvector v2 (lambda=-1, flipping)
  \draw[->, thick, red!70!black] (0,0) -- (1,-0.5)
    node[below right] {$\vec{v}_2$};
  \draw[->, thick, red, dashed] (0,0) -- (-1,0.5)
    node[above left] {$f(\vec{v}_2)=-\vec{v}_2$};

  % Generic vector
  \draw[->, thick, black!50!green] (0,0) -- (1.5,0.5)
    node[right] {$\vec{w}$};
  \draw[->, thick, black!50!green, dashed] (0,0) -- (2.5,1.5)
    node[right] {$f(\vec{w})$};

  % Legend
  \node[anchor=north west, text width=5cm] at (-2.4,3.5) {%
    \small\textcolor{blue!70!black}{$\lambda_1=2$~: étirement} \\
    \small\textcolor{red!70!black}{$\lambda_2=-1$~: retournement}};
\end{tikzpicture}
\caption{Action d'un endomorphisme sur ses vecteurs propres~:
  l'endomorphisme étire le long de $\vec{v}_1$ (valeur propre $2$) et
  retourne le long de $\vec{v}_2$ (valeur propre $-1$).}
\label{fig:vp-illustration}
\end{figure}

\begin{figure}[ht]
\centering
\begin{tikzpicture}[>=Stealth, node distance=2.5cm]
  % Nodes
  \node (A) {$A\in\Mat_n(\K)$};
  \node[right=3.5cm of A] (D) {$D=\diag(\lambda_1,\ldots,\lambda_n)$};
  \node[below=1.8cm of A] (f) {$f\in\End(E)$};
  \node[below=1.8cm of D] (fdiag)
    {$f$ diagonal dans $\cB'$};

  % Arrows
  \draw[->] (A) -- node[above] {$\inv{P}(\cdot)P$} (D);
  \draw[<->] (A) -- node[left] {$\Mat_\cE$} (f);
  \draw[<->] (D) -- node[right] {$\Mat_{\cB'}$} (fdiag);
  \draw[->] (f) -- node[below] {changement de base} (fdiag);

  % Extra annotation
  \node[above=0.3cm of A, text=blue!60!black] {\small base $\cE$};
  \node[above=0.3cm of D, text=red!60!black]
    {\small base $\cB'$ de vecteurs propres};
\end{tikzpicture}
\caption{Diagramme de diagonalisation~: le passage de la matrice $A$
  (dans la base canonique $\cE$) à la matrice diagonale $D$ (dans la
  base de vecteurs propres $\cB'$) via la matrice de passage $P$.}
\label{fig:diag-diagram}
\end{figure}

% ============================================================
\section{Exercices}\label{sec:exos-vp}
% ============================================================

\begin{exercise}{\basic{} Valeurs propres d'une matrice $2\times 2$}{exo:vp-2x2}
  Soit $A=\begin{pmatrix}5&-3\\6&-4\end{pmatrix}$.
  \begin{enumerate}[label=(\alph*)]
    \item Calculer le polynôme caractéristique de $A$.
    \item Déterminer les valeurs propres et les sous-espaces propres.
    \item La matrice $A$ est-elle diagonalisable~? Si oui, la
      diagonaliser.
  \end{enumerate}
\end{exercise}

\begin{exercise}{\basic{} Valeurs propres d'une matrice $3\times 3$}{exo:vp-3x3}
  Soit $A=\begin{pmatrix}2&0&0\\0&3&-1\\0&-1&3\end{pmatrix}$.
  \begin{enumerate}[label=(\alph*)]
    \item Calculer $\chi_A(\lambda)$.
    \item Déterminer $\Spec(A)$ et les sous-espaces propres.
    \item Diagonaliser $A$.
  \end{enumerate}
\end{exercise}

\begin{exercise}{\basic{} Polynôme caractéristique et trace/déterminant}{exo:chi-trace-det}
  Soit $A=\begin{pmatrix}a&b\\c&d\end{pmatrix}\in\Mat_2(\R)$.
  \begin{enumerate}[label=(\alph*)]
    \item Montrer que $\chi_A(\lambda)=\lambda^2-\tr(A)\lambda+\det(A)$.
    \item En déduire que si $\lambda_1,\lambda_2$ sont les valeurs
      propres de $A$ (dans $\C$), alors $\lambda_1+\lambda_2=\tr(A)$
      et $\lambda_1\lambda_2=\det(A)$.
    \item Généraliser au cas $n\times n$~: montrer que la somme des
      valeurs propres (comptées avec multiplicité) vaut $\tr(A)$ et
      leur produit vaut $\det(A)$.
  \end{enumerate}
\end{exercise}

\begin{exercise}{\basic{} Vérification de Cayley-Hamilton}{exo:cayley-hamilton-verif}
  Soit $A=\begin{pmatrix}2&1\\0&3\end{pmatrix}$.
  \begin{enumerate}[label=(\alph*)]
    \item Calculer $\chi_A(\lambda)$.
    \item Calculer $A^2$, puis vérifier que $\chi_A(A)=0$.
    \item En déduire une expression de $\inv{A}$ comme polynôme en $A$.
  \end{enumerate}
\end{exercise}

\begin{exercise}{\intermediate{} Matrice non diagonalisable}{exo:non-diag}
  Soit $A=\begin{pmatrix}2&1\\0&2\end{pmatrix}$.
  \begin{enumerate}[label=(\alph*)]
    \item Déterminer $\chi_A(\lambda)$.
    \item Trouver les valeurs propres et les sous-espaces propres.
    \item Comparer multiplicités algébrique et géométrique. Conclure
      sur la diagonalisabilité de $A$.
    \item Calculer néanmoins $A^n$ pour tout $n\in\N$ en écrivant
      $A=2I_2+N$ avec $N=\begin{pmatrix}0&1\\0&0\end{pmatrix}$ et
      en utilisant la formule du binôme.
  \end{enumerate}
\end{exercise}

\begin{exercise}{\intermediate{} Puissance d'une matrice}{exo:puissance-mat}
  Soit $A=\begin{pmatrix}3&-2\\1&0\end{pmatrix}$.
  \begin{enumerate}[label=(\alph*)]
    \item Diagonaliser $A$.
    \item En déduire $A^n$ pour tout $n\in\N$.
    \item Calculer $A^{10}$.
  \end{enumerate}
\end{exercise}

\begin{exercise}{\intermediate{} Récurrence linéaire}{exo:recurrence}
  On considère la suite définie par $u_0=1$, $u_1=2$ et
  $u_{n+2}=5u_{n+1}-6u_n$.
  \begin{enumerate}[label=(\alph*)]
    \item Écrire cette récurrence sous forme matricielle~:
      $X_{n+1}=AX_n$ avec $X_n=\begin{pmatrix}u_{n+1}\\u_n\end{pmatrix}$.
    \item Diagonaliser la matrice $A$.
    \item En déduire l'expression de $u_n$ en fonction de $n$.
  \end{enumerate}
\end{exercise}

\begin{exercise}{\intermediate{} Polynôme minimal}{exo:poly-min}
  Pour chacune des matrices suivantes, déterminer le polynôme minimal
  et dire si la matrice est diagonalisable~:
  \begin{enumerate}[label=(\alph*)]
    \item $A=\begin{pmatrix}3&0\\0&3\end{pmatrix}$.
    \item $B=\begin{pmatrix}3&1\\0&3\end{pmatrix}$.
    \item $C=\begin{pmatrix}1&0&0\\0&2&0\\0&0&1\end{pmatrix}$.
    \item $D=\begin{pmatrix}1&0&0\\0&1&0\\0&0&2\end{pmatrix}$.
  \end{enumerate}
\end{exercise}

\begin{exercise}{\intermediate{} Diagonalisation et système différentiel}{exo:syst-diff}
  Résoudre le système
  $\begin{cases}x'(t)=3x(t)-y(t)\\y'(t)=-x(t)+3y(t)\end{cases}$
  en diagonalisant la matrice du système.
\end{exercise}

\begin{exercise}{\intermediate{} Chaîne de Markov}{exo:markov}
  Dans une ville, chaque jour, $80\%$ des usagers du bus reprennent le
  bus et $20\%$ passent au vélo~; $40\%$ des cyclistes reprennent le
  vélo et $60\%$ passent au bus.
  \begin{enumerate}[label=(\alph*)]
    \item Écrire la matrice de transition $M$.
    \item Diagonaliser $M$.
    \item Déterminer le vecteur de probabilités stationnaires.
    \item Si initialement tout le monde prend le bus, quelle est la
      répartition après $n$ jours~? Trouver la limite quand
      $n\to+\infty$.
  \end{enumerate}
\end{exercise}

\begin{exercise}{\challenging{} Diagonalisation simultanée}{exo:diag-simult}
  Soient $A,B\in\Mat_n(\K)$ deux matrices diagonalisables telles que
  $AB=BA$.
  \begin{enumerate}[label=(\alph*)]
    \item Montrer que pour toute valeur propre $\lambda$ de $A$, le
      sous-espace propre $E_\lambda(A)$ est stable par $B$.
    \item En déduire que $A$ et $B$ sont \emph{simultanément
      diagonalisables}~: il existe $P\in\GL_n(\K)$ telle que
      $\inv{P}AP$ et $\inv{P}BP$ soient toutes deux diagonales.
    \item Application~: diagonaliser simultanément
      $A=\begin{pmatrix}1&2\\2&1\end{pmatrix}$ et
      $B=\begin{pmatrix}3&-2\\-2&3\end{pmatrix}$.
  \end{enumerate}
\end{exercise}

\begin{exercise}{\challenging{} Endomorphisme et polynôme annulateur}{exo:annulateur}
  Soit $f\in\End(E)$ avec $\dim E=n$. On suppose que $f^2=f$ (i.e.,
  $f$ est un projecteur).
  \begin{enumerate}[label=(\alph*)]
    \item Montrer que les seules valeurs propres possibles de $f$
      sont $0$ et $1$.

      \emph{Indication~:} si $f(\vec{v})=\lambda\vec{v}$, calculer
      $f^2(\vec{v})$.
    \item Montrer que $\mu_f$ divise $X^2-X=X(X-1)$.
    \item En déduire que $f$ est diagonalisable.
    \item Montrer que $E=\Ker f\oplus\im f$ et interpréter les
      sous-espaces propres.
  \end{enumerate}
\end{exercise}

\begin{exercise}{\challenging{} Cayley-Hamilton et calcul d'inverse}{exo:cayley-inv}
  Soit $A=\begin{pmatrix}1&1&1\\0&2&1\\0&0&3\end{pmatrix}$.
  \begin{enumerate}[label=(\alph*)]
    \item Calculer $\chi_A(\lambda)$.
    \item En utilisant le théorème de Cayley-Hamilton, exprimer
      $\inv{A}$ comme un polynôme en $A$ (de degré $\leq 2$).
    \item Calculer $A^{-1}$ par cette méthode et vérifier le résultat.
  \end{enumerate}
\end{exercise}

\begin{exercise}{\challenging{} Trigonalisation}{exo:trigonalisation}
  Soit $A=\begin{pmatrix}2&1&0\\0&2&1\\0&0&2\end{pmatrix}$.
  \begin{enumerate}[label=(\alph*)]
    \item Montrer que $A$ n'est pas diagonalisable.
    \item Déterminer le polynôme minimal de $A$.
    \item Calculer $A^n$ en écrivant $A=2I_3+N$ où $N$ est nilpotente.
      Déterminer l'indice de nilpotence de $N$ et utiliser la formule
      du binôme $(2I_3+N)^n=\sum_{k=0}^{2}\binom{n}{k}2^{n-k}N^k$.
  \end{enumerate}
\end{exercise}

\begin{exercise}{\challenging{} Spectre et trace}{exo:spectre-trace}
  Soit $A\in\Mat_n(\K)$ de valeurs propres
  $\lambda_1,\ldots,\lambda_n$ (comptées avec multiplicité, dans une
  clôture algébrique).
  \begin{enumerate}[label=(\alph*)]
    \item Montrer que $\tr(A^k)=\sum_{i=1}^n\lambda_i^k$ pour tout
      $k\geq 1$.

      \emph{Indication~:} trigonaliser $A$ sur la clôture algébrique.
    \item En déduire que si $A^2=A$, alors $\tr(A)=\rg(A)$.
    \item Montrer que si $A$ est nilpotente, alors $\tr(A^k)=0$ pour
      tout $k\geq 1$.
  \end{enumerate}
\end{exercise}

% ============================================================
\section*{Résumé du chapitre}
\addcontentsline{toc}{section}{Résumé du chapitre}
% ============================================================

\begin{framed}
\textbf{Ce qu'il faut retenir.}

\begin{enumerate}[label=\textbf{\arabic*.},leftmargin=2em]
  \item \textbf{Valeur propre / vecteur propre.} $\lambda\in\K$ est
    valeur propre de $f$ si $\exists\,\vec{v}\neq\vzero$ tel que
    $f(\vec{v})=\lambda\vec{v}$, i.e.\
    $\Ker(f-\lambda\Id)\neq\set{\vzero}$.

  \item \textbf{Sous-espace propre.}
    $E_\lambda(f)=\Ker(f-\lambda\Id)$ est un sous-espace vectoriel.
    Les sous-espaces propres associés à des valeurs propres distinctes
    sont en somme directe.

  \item \textbf{Polynôme caractéristique.}
    $\chi_f(\lambda)=\det(f-\lambda\Id)$, de degré $n$, avec
    $\chi_f(0)=\det(f)$ et coefficient de $(-1)^{n-1}\lambda^{n-1}$
    égal à $\tr(f)$. Les valeurs propres sont les racines de $\chi_f$.

  \item \textbf{Cayley-Hamilton.} $\chi_f(f)=0$ --- toute matrice est
    racine de son polynôme caractéristique.

  \item \textbf{Multiplicités.} Pour toute valeur propre~$\lambda$~:
    $1\leq m_g(\lambda)\leq m_a(\lambda)$.

  \item \textbf{Diagonalisabilité.} $f$ est diagonalisable ssi $\chi_f$
    est scindé et $m_g=m_a$ pour chaque valeur propre, ssi $\mu_f$ est
    scindé à racines simples.

  \item \textbf{Trigonalisation.} $f$ est trigonalisable ssi $\chi_f$
    est scindé. Sur $\C$, tout endomorphisme est trigonalisable.

  \item \textbf{Polynôme minimal.} $\mu_f$ est le polynôme unitaire
    annulateur de plus petit degré. Il divise $\chi_f$ et a les mêmes
    racines.

  \item \textbf{Applications.} Si $A=PD\inv{P}$~:
    \begin{itemize}
      \item $A^k=PD^k\inv{P}$ (puissances).
      \item Résolution de récurrences linéaires ($u_{n+k}=\sum a_i u_{n+i}$).
      \item Chaînes de Markov~: l'état stationnaire est le vecteur
        propre pour $\lambda=1$.
      \item Systèmes différentiels~: $X'=AX$ se résout par
        diagonalisation.
    \end{itemize}

  \item \textbf{Formules clés.}
    \begin{itemize}
      \item $\sum\lambda_i=\tr(A)$, \;$\prod\lambda_i=\det(A)$.
      \item $A$ inversible $\Leftrightarrow$ $0\notin\Spec(A)$.
      \item Matrices semblables $\Rightarrow$ même $\chi$, même $\mu$,
        même spectre, même trace, même déterminant.
    \end{itemize}
\end{enumerate}
\end{framed}